\section{Related work} \label{sec:realted}

\paragraph{Multi-turn reasoning degradation.}
Recent work documents systematic performance loss in multi-turn settings. \citet{lanban2025lost} report an average 39\% accuracy drop
across six generation tasks over 200{,}000 simulated conversations with 15 frontier models, characterizing failure as a ``lost-in-conversation''
phenomenon. \citet{liu2024lost} show that language models do not robustly use information in long contexts, with retrieval accuracy varying 
sharply by position. These findings establish multi-turn degradation as broad-spectrum, but existing multi-turn benchmarks largely target 
underspecification and clarification rather than numerical chaining. Our dataset complements this literature by isolating numerical-dependency
failures, where each sub-question is fully specified but its solution requires the prior turn's numerical result.

\paragraph{Prompt and token-level compression}
Prompt compression methods reduce input length while attempting to preserve task performance. LLMLingua~\citep{jiang2023llmlingua} and its 
extensions identify and remove low-information tokens via small auxiliary models, operating at the lexical level without awareness of which 
tokens encode task-critical structure. By design, a numerical constant inside an inequality and a discourse adjective receive comparable
treatment. SLSC differs along two axes: compression operates over semantic units---atomic mathematical propositions---rather than tokens, and 
is performed by a dedicated curator agent prompted to identify mathematical content. We view the two approaches as orthogonal and do not include 
token-level compression in our experimental comparison, which focuses on state-maintenance strategies sharing the same multi-turn decomposed 
setup.

\paragraph{Reasoning-state compression and summarization}
A closer line of work compresses intermediate reasoning state rather than raw inputs. Reflexion~\citep{shinn2023reflexion} condenses prior
trajectories into natural-language reflections that condition subsequent attempts. InftyThink~\citep{yan2025inftythink} extends this idea to
long-horizon reasoning by iteratively  summarizing chain-of-thought and re-prompt the model on the summary, achieving large context savings.
Conceptually these methods sit closest to SLSC, and two differences are central. First, both produce flat natural-language summaries in which
numerical constraints stated in passing receive no privileged representation; SLSC enforces an atomic-proposition schema in which each targets
monolithic single-question reasoning with internally-induced summarization points and assumes a trained summarization checkpoint, whereas
our setting is externally-decomposed multi-turn dialogue with explicit numerical dependencies and is fully training-free. A  head-to-head
empirical comparison would require either retraining InftyThink on our dependency-annotated dataset or rebuilding it as a prompted pipeline;
we leave this to future work and restrict our experiments to training-free, setup-compatible baselines.

\paragraph{Implicit assumptions in compression literature}
The compression methods above share an implicit assumption: that reducing context length improves downstream task performance, either by
mitigating attention dilution~\citep{liu2024lost} or by enabling longer reasoning chains within fixed context budgets. This assumption is
rarely tested in controlled settings where compressed and uncompressed strategies are compared on identical problems with statistical rigor.
Our work directly evaluates this assumption in the multi-turn mathematical reasoning regime. We find that compression methods substantially
outperform raw history accumulation in accuracy at 4B parameters and above (Section~\ref{sec:results}), while delivering substantial per-turn
token savings. This positions compression as both an accuracy intervention and a serving-cost intervention, and suggests that prior claims of
compression-driven accuracy gains are well-founded in this regime, though the benefit comes primarily from bounded context size rather than
from structured representation per se. 

\paragraph{Formal language augmentation for reasoning.}
A separate line of work explores formal-language augmentation for mathematical reasoning. Autoformalization~\citep{wu2022autoformalization}
translates natural-language statements into formal expressions; Draft, Sketch, and Prove~\citep{jiang2023draft} interleaves natural-language
drafting with formal verification; LeanDojo~\citep{yang2023leandojo} provides infrastructure for retrieval-augmented theorem proving. These
works target end-to-end proof construction or single-statement formalization, with verifier feedback in the loop. SLSC's Lean~4 augmentation
(Section~\ref{sec:agent_b}) instead explores a lighter use of formal language---type-annotating extracted propositions without a verifier---
and our ablation (Section~\ref{sec:lean_ablation}) finds consistent accuracy degradation. This negative result aligns with the broader emphasis on
verifier-grounded formalization: absent a Lean kernel check, the type annotation introduces noise that the reasoner cannot reliably exploit.
SLSC thus occupies a specific intersection in this landscape---training-free, multi-turn, mathematics-targeted, and operating on semantic units
of intermediate state---with raw context accumulation and natural-language summary as the primary setup-compatible baselines.